\documentclass[../../main.tex]{subfiles}
\begin{document}

平面上に点列 $P_0,P_1,P_2,\cdots,P_n,\cdots$ があり，
$P_0$，$P_1$ の座標はそれぞれ $(0,0)$，$(1,0)$ である．
また，任意の自然数 $n$ に対し，
線分 $P_nP_{n+1}$ の長さは，線分 $P_{n-1}P_n$ の長さの2倍で，
半直線 $P_nP_{n+1}$ が半直線 $P_{n-1}P_n$ となす角は $120^\circ$ である．
$P_{3n}$ の座標を求めよ．

\begin{figure}[htbp]
  \centering
  \begin{tikzpicture}[scale=1.5, >=stealth]
    \draw[->] (-1.8,0) -- (2.0,0) node[right] {$x$};
    \draw[->] (0,-3.2) -- (0,2.5) node[above] {$y$};

    \coordinate (P0) at (0,0);
    \coordinate (P1) at (1,0);
    \coordinate (P2) at (0, 1.732);
    \coordinate (P3) at (-1.5, -2.598);

    \draw[thick] (P0) -- (P1) -- (P2) -- (P3);

    \node[circle, fill, inner sep=1.5pt] at (P0) {}; \node[below left] at (P0) {$P_0$};
    \node[circle, fill, inner sep=1.5pt] at (P1) {}; \node[below right] at (P1) {$P_1$};
    \node[circle, fill, inner sep=1.5pt] at (P2) {}; \node[above right] at (P2) {$P_2$};
    \node[circle, fill, inner sep=1.5pt] at (P3) {}; \node[above left] at (P3) {$P_3$};

    \draw[->] (P1) +(0:0.4) arc (0:120:0.4);
    \node at ($(P1)+(60:0.55)$) {$120^\circ$};

    \draw[->] (P2) +(120:0.4) arc (120:240:0.4);
    \node at ($(P2)+(180:0.65)$) {$120^\circ$};

    \draw[->] (P3) +(240:0.4) arc (240:360:0.4);
    \node at ($(P3)+(-60:0.55)$) {$120^\circ$};
  \end{tikzpicture}
\end{figure}

\end{document}
